% arara: lualatex_git
\documentclass[4x6]{recipe-card}

\begin{document}

% --- FRONT SIDE ---
\begin{recipe}{Negroni}
\setBaseSpirit{Gin}
\noindent\begin{minipage}{\textwidth}
\small\raggedright
\cocktailGlassWithLabel{Rocks}\quad
\cocktailMethodWithLabel{Stir}\quad
\end{minipage}\par\vspace{6pt}
\begin{ingredients}
\ingredient{1 oz}{Gin}
\ingredient{1 oz}{Campari}
\ingredient{1 oz}{Sweet Vermouth}
\group{Garnish}
\ingredient{}{Orange Peel}
\group{Glassware}
\ingredient{}{Rocks}
\group{Preparation}
\ingredient{}{Stir}
\end{ingredients}
\begin{process}
\step{Add ingredients to mixing glass with ice.}
\step{Stir until chilled.}
\step{Strain into rocks glass.}
\end{process}

\vfill
\begin{recipeinfo}
    Classic equal-parts cocktail.
\end{recipeinfo}

% --- BACK SIDE ---
\back
\drawBackMirrorSpiritBar
\begin{center}
    \footnotesize\textsc{\color{recipeRed} History \& Provenance}
    \vspace{-3pt}
    \rule{0.92\textwidth}{0.35pt}
\end{center}
\vspace{-5pt}
\begin{multicols*}{2}
\scriptsize
\setlength{\columnsep}{9pt}
The Negroni is widely traced to early 20th-century Florence, where Count Camillo Negroni asked for his Americano to be strengthened with gin instead of soda.
\par
This simple substitution created a bold, bitter cocktail that became a cornerstone of aperitivo culture.
\par
Today, the Negroni remains one of the most influential classic cocktails, inspiring countless variations.
\vspace{2pt}
\begin{flushright}
\qrcode[height=0.72in]{https://poly-repo.github.io/recipes/cocktails/negroni}
\end{flushright}
\end{multicols*}

\end{recipe}

% --- FRONT SIDE ---
\begin{recipe}{Negroni}
\setBaseSpirit{Gin}
\noindent\begin{minipage}{\textwidth}
\small\raggedright
\cocktailGlassWithLabel{Rocks}\quad
\cocktailMethodWithLabel{Stir}\quad
\end{minipage}\par\vspace{6pt}
\begin{ingredients}
\ingredient{1 oz}{Gin}
\ingredient{1 oz}{Campari}
\ingredient{1 oz}{Sweet Vermouth}
\group{Garnish}
\ingredient{}{Orange Peel}
\group{Glassware}
\ingredient{}{Rocks}
\group{Preparation}
\ingredient{}{Stir}
\end{ingredients}
\begin{process}
\step{Add ingredients to mixing glass with ice.}
\step{Stir until chilled.}
\step{Strain into rocks glass.}
\end{process}

\vfill
\begin{recipeinfo}
    Classic equal-parts cocktail.
\end{recipeinfo}

% --- BACK SIDE ---
\back
\drawBackMirrorSpiritBar
\begin{center}
    \footnotesize\textsc{\color{recipeRed} History \& Provenance}
    \vspace{-3pt}
    \rule{0.92\textwidth}{0.35pt}
\end{center}
\vspace{-5pt}
\begin{multicols*}{2}
\scriptsize
\setlength{\columnsep}{9pt}
The Negroni is widely traced to early 20th-century Florence, where Count Camillo Negroni asked for his Americano to be strengthened with gin instead of soda.
\par
This simple substitution created a bold, bitter cocktail that became a cornerstone of aperitivo culture.
\par
Today, the Negroni remains one of the most influential classic cocktails, inspiring countless variations.
\vspace{2pt}
\begin{flushright}
\qrcode[height=0.72in]{https://poly-repo.github.io/recipes/cocktails/negroni}
\end{flushright}
\end{multicols*}

\end{recipe}

\end{document}
